\Article{30}

When the modulus is composite, the following method is sometimes preferable.

Let the modulus be $=mn$ and the given congruence $ax\equiv b$. First solve this congruence modulo $m$, and suppose that it is satisfied when $x\equiv v \ (\mathrm{mod.}\ \frac{m}{\delta})$, where $\delta$ denotes the greatest common divisor of $m,a$. Every value of $x$ satisfying the congruence $ax\equiv b$ modulo $mn$ clearly also satisfies it modulo $m$, and so has the form $v+\frac{m}{\delta}x'$, where $x'$ is an indeterminate number. Conversely, however, not every number of the form $v+\frac{m}{\delta}x'$ satisfies the congruence modulo $mn$. To determine $x'$ so that $v+\frac{m}{\delta}x'$ is a root of the congruence $ax\equiv b \ (\mathrm{mod.}\ mn)$, we can solve the congruence $\frac{am}{\delta}x'+av\equiv b\ (\mathrm{mod.}\ mn)$, which is equivalent to $\frac{a}{\delta}x'\equiv\frac{b-av}{m}\ (\mathrm{mod.}\ n)$. Thus the solution of any congruence of the first degree modulo $mn$ can be reduced to the solution of two congruences modulo $m$ and $n$. If $n$ is itself a product of two factors, the solution of the congruence modulo $n$ similarly depends on solving two congruences whose moduli are those factors. In general, solving a congruence with any composite modulus depends on solving other congruences whose moduli are factors of that number. These factors can always be chosen to be primes, if that seems convenient.

\emph{Example.} Consider the congruence $19x\equiv 1 \ (\mathrm{mod.}\ 140)$. First solve it modulo $2$, giving $x\equiv 1\ (\mathrm{mod.}\ 2)$. Put $x=1+2x'$; then $38x'\equiv -18\ (\mathrm{mod.}\ 140)$, which is equivalent to $19x'\equiv -9\ (\mathrm{mod.}\ 70)$. If this is
\SourceLeaf{28}
again solved modulo $2$, we get $x'\equiv 1\ (\mathrm{mod.}\ 2)$. Setting $x'=1+2x''$ gives $38x''\equiv -28\ (\mathrm{mod.}\ 70)$, or $19x''\equiv -14\ (\mathrm{mod.}\ 35)$. Solving this modulo $5$ gives $x''\equiv 4\ (\mathrm{mod.}\ 5)$, and substituting $x''=4+5x'''$ gives $95x'''\equiv -90\ (\mathrm{mod.}\ 35)$, or $19x'''\equiv -18\ (\mathrm{mod.}\ 7)$. Finally, this gives $x'''\equiv 2\ (\mathrm{mod.}\ 7)$, and on setting $x'''=2+7x''''$ we obtain $x=59+140x''''$. So $x\equiv 59\ (\mathrm{mod.}\ 140)$ is the complete solution of the given congruence.

\Article{31}

Just as the root of the equation $ax=b$ is expressed by $\frac{b}{a}$, we will denote any root of the congruence $ax\equiv b$ by $\frac{b}{a}$, adding the modulus to distinguish the congruence. For example, $\frac{19}{17}\ (\mathrm{mod.}\ 12)$ denotes any number that is $\equiv 11\ (\mathrm{mod.}\ 12)$\begingroup\renewcommand{\thefootnote}{*)}\footnote{which, by analogy, can be denoted by $\frac{11}{1}\ (\mathrm{mod.}\ 12)$.}\endgroup. In general, what precedes shows that $\frac{a}{b}\ (\mathrm{mod.}\ c)$ represents nothing real (or, if one prefers, something imaginary) if $a,c$ have a common divisor that does not divide $b$. Except in this case, the expression $\frac{b}{a}\ (\mathrm{mod.}\ c)$ always has real values, and indeed infinitely many. These are all congruent modulo $c$ when $a$ is coprime to $c$, or modulo $\frac{c}{\delta}$ when $\delta$ is the greatest common divisor of $c,a$.

The calculation rules for these expressions are much like those for ordinary fractions. Here are a few properties that follow easily from what precedes.

\hangindent=1.7em\hangafter=1
1.\quad If, modulo $c$, $a\equiv\alpha$, $b\equiv\beta$, then the expressions $\frac{a}{b}\ (\mathrm{mod.}\ c)$ and $\frac{\alpha}{\beta}\ (\mathrm{mod.}\ c)$ are equivalent.

\hangindent=1.7em\hangafter=1
2.\quad $\frac{a\delta}{b\delta}\ (\mathrm{mod.}\ c\delta)$ and $\frac{a}{b}\ (\mathrm{mod.}\ c)$ are equivalent.

\hangindent=1.7em\hangafter=1
3.\quad $\frac{ak}{bk}\ (\mathrm{mod.}\ c)$ and $\frac{a}{b}\ (\mathrm{mod.}\ c)$ are equivalent when $k$ is coprime to $c$.

Many other similar propositions could be given. But they present no difficulty and are not especially needed for what follows, so we move on.

\Needspace{6\baselineskip}
\begin{center}\small\emph{Finding a number congruent to given residues modulo given numbers.}\end{center}

\Article{32}

The problem of \emph{finding all numbers that give specified residues modulo any number of given moduli} will be very useful later, and is easily solved using what precedes. First let there be two moduli $A,B$, for which the number
\SourceLeaf{29}
we seek, $z$, is to be congruent to $a,b$, respectively. All values of $z$ therefore have the form $Ax+a$, where $x$ is indeterminate but must satisfy $Ax+a\equiv b\ (\mathrm{mod.}\ B)$. If the greatest common divisor of $A,B$ is $\delta$, the complete solution of this congruence has the form $x\equiv v\ (\mathrm{mod.}\ \frac{B}{\delta})$, or equivalently $x=v+\frac{kB}{\delta}$, where $k$ is an arbitrary integer. The formula $Av+a+\frac{kAB}{\delta}$ therefore includes all values of $z$, that is, $z\equiv Av+a\ (\mathrm{mod.}\ \frac{AB}{\delta})$ is the complete solution of the problem. If the moduli $A,B$ are joined by a third, $C$, for which the desired number $z$ is to be $\equiv c$, we clearly proceed in the same way, since the first two conditions have already been combined into one. Specifically, if the greatest common divisor of $\frac{AB}{\delta}.$ and $C$ is $=\varepsilon$, and the congruence $\frac{AB}{\delta}x+Av+a\equiv c\ (\mathrm{mod.}\ C)$ has solution $x\equiv w\ (\mathrm{mod.}\ \frac{C}{\varepsilon})$, then the congruence $z\equiv\frac{ABw}{\delta}+Av+a\ (\mathrm{mod.}\ \frac{ABC}{\delta\varepsilon})$ solves the problem completely. We proceed similarly however many moduli are given. Notice that $\frac{AB}{\delta}$, $\frac{ABC}{\delta\varepsilon}$ are the least common multiples of $A,B$ and $A,B,C$, respectively. So, for any number of moduli $A,B,C$ etc., if their least common multiple is $M$, the complete solution has the form $z\equiv r\ (\mathrm{mod.}\ M)$. If any of the auxiliary congruences is unsolvable, we must conclude that the problem is impossible. Clearly this cannot happen when all the numbers $A,B,C$ etc. are pairwise coprime.

\emph{Example.} Let $A,B,C$ and $a,b,c$ be $504,35,16$ and $17,-4,33$. The two conditions that $z$ be $\equiv 17\ (\mathrm{mod.}\ 504)$ and $\equiv -4\ (\mathrm{mod.}\ 35)$ are equivalent to the single condition that it be $\equiv 521\ \mathrm{mod.}\ 2520)$. Combining this with $z\equiv 33\ (\mathrm{mod.}\ 16)$ gives $z\equiv 3041\ (\mathrm{mod.}\ 5040)$.

\Article{33}

When the numbers $A,B,C$ etc. are pairwise coprime, their product is known to be their least common multiple. In that case, all the congruences $z\equiv a\ (\mathrm{mod.}\ A)$, $z\equiv b\ (\mathrm{mod.}\ B)$ etc. are clearly equivalent to the single congruence $z\equiv r\ (\mathrm{mod.}\ R)$, where $R$ is the product of $A,B,C$ etc. Conversely, the single condition $z\equiv r\ (\mathrm{mod.}\ R)$ can therefore be split into several. If $R$ is factored in any way into pairwise coprime factors $A,B,C$ etc., the conditions $z\equiv r\ (\mathrm{mod.}\ A)$, $z\equiv r\ (\mathrm{mod.}\ B)$, $z\equiv r\ (\mathrm{mod.}\ C)$ etc. meet the requirement in full. This observation gives us a way
\SourceLeaf{30}
not only to detect immediately any impossibility in the given conditions, but also to make the calculation more convenient and orderly.

\Article{34}

As above, let the given conditions be $z\equiv a\ (\mathrm{mod.}\ A)$, $z\equiv b\ (\mathrm{mod.}\ B)$, $z\equiv c\ (\mathrm{mod.}\ C)$. Factor each modulus into pairwise coprime factors: $A$ into $A'A''A'''$ etc., $B$ into $B'B''B'''$ etc., and so on. Choose these so that the numbers $A',A''$ etc., $B',B''$ etc. are primes or powers of primes. If one of the numbers $A,B,C$ etc. is already a prime or a power of a prime, it needs no factorization. What precedes shows that we can replace the given conditions by $z\equiv a\ (\mathrm{mod.}\ A')$, $z\equiv a\ (\mathrm{mod.}\ A'')$, $z\equiv a\ (\mathrm{mod.}\ A''')$ etc., $z\equiv b\ (\mathrm{mod.}\ B')$, $z\equiv b\ (\mathrm{mod.}\ B'')$ etc. Now, unless all the numbers $A,B,C$ etc. are pairwise coprime, for example if $A$ is not coprime to $B$, the prime divisors of $A,B$ clearly cannot all be different. Among the factors $A',A'',A'''$ etc., at least one must have an equal, a multiple or a divisor among $B',B'',B'''$ etc. \emph{First}, if $A'=B'$, the conditions $z\equiv a\ (\mathrm{mod.}\ A')$, $z\equiv b\ (\mathrm{mod.}\ B')$ must be identical, that is, $a\equiv b\ (\mathrm{mod.}\ A'\ \text{or}\ B')$, so either can be discarded. If $a\equiv b\ (\mathrm{mod.}\ A')$ does not hold, the problem is impossible. \emph{Second}, if $B'$ is a multiple of $A'$, the condition $z\equiv a\ (\mathrm{mod.}\ A')$ must be included in $z\equiv b\ (\mathrm{mod.}\ B')$. In other words, the condition $z\equiv b\ (\mathrm{mod.}\ A')$, which follows from the latter, must be identical to the former. So the condition $z\equiv a\ (\mathrm{mod.}\ A')$ can be discarded unless it conflicts with the other condition, in which case the problem is impossible. After all redundant conditions have been discarded in this way, all the remaining moduli among $A',A'',A'''$ etc., $B',B'',B'''$ etc. are clearly pairwise coprime. We can then be sure that the problem is possible and proceed by the rules already given.

\Article{35}

\emph{Example.} If, as above, we require $z\equiv 17\ (\mathrm{mod.}\ 504)$, $\equiv -4\ (\mathrm{mod.}\ 35)$ and $\equiv 33\ (\mathrm{mod.}\ 16)$, these conditions can be split into $z\equiv 17\ (\mathrm{mod.}\ 8)$, $\equiv 17\ (\mathrm{mod.}\ 9)$, $\equiv 17\ (\mathrm{mod.}\ 7)$, $\equiv -4\ (\mathrm{mod.}\ 5)$, $\equiv -4\ (\mathrm{mod.}\ 7)$, $\equiv 33\ (\mathrm{mod.}\ 16)$. We can discard $z\equiv 17\ (\mathrm{mod.}\ 8)$, $z\equiv 17\ (\mathrm{mod.}\ 7)$, since the former is included in $z\equiv 33\ (\mathrm{mod.}\ 16)$, and the latter is identical to $z\equiv -4\ (\mathrm{mod.}\ 7)$. This leaves

\SourceLeaf{31}\[
z\equiv\left\{\begin{array}{rl}
17 & (\mathrm{mod.}\ 9)\\
-4 & (\mathrm{mod.}\ 5)\\
-4 & (\mathrm{mod.}\ 7)\\
33 & (\mathrm{mod.}\ 16)
\end{array}\right.
\quad\text{so}\quad z\equiv 3041\ (\mathrm{mod.}\ 5040).
\]

It is usually more convenient to recombine separately those remaining conditions that came from the same original condition, since this is easily done. For \emph{example}, if some of the conditions $z\equiv a\ (\mathrm{mod.}\ A')$, $z\equiv a\ (\mathrm{mod.}\ A'')$ etc. have been discarded, the condition reconstructed from the rest is $z\equiv a$ modulo the product of all the remaining moduli among $A',A'',A'''$ etc. Thus, in our example, the conditions $z\equiv -4\ (\mathrm{mod.}\ 5)$, $z\equiv -4\ (\mathrm{mod.}\ 7)$ immediately give back the condition $z\equiv -4\ (\mathrm{mod.}\ 35)$ from which they came. This also shows that, as far as shortening the calculation is concerned, it does matter which redundant conditions we discard. But explaining this and other practical devices, which are learned much more easily through use than through rules, is not our purpose here.

\Article{36}

When all the moduli $A,B,C,D$ etc. are pairwise coprime, the following method is often preferable. Find a number $\alpha$ that is congruent to unity modulo $A$ and to zero modulo the product of the other moduli. Equivalently, let $\alpha$ be any value (usually the \emph{least}) of the expression $\frac{1}{BCD\ \mathrm{etc.}}\ (\mathrm{mod.}\ A)$, multiplied by $BCD$ etc. (see Article 32). Similarly, let $\beta\equiv 1\ (\mathrm{mod.}\ B)$ and $\equiv 0\ (\mathrm{mod.}\ ACD\ \mathrm{etc.})$, $\gamma\equiv 1\ (\mathrm{mod.}\ C)$ and $\equiv 0\ (\mathrm{mod.}\ ABD\ \mathrm{etc.})$, etc. Then, if we seek a number $z$ congruent modulo $A,B,C,D$ etc. to $a,b,c,d$ etc., respectively, we can put

\[
z\equiv\alpha a+\beta b+\gamma c+\delta d\ \mathrm{etc.}\ (\mathrm{mod.}\ ABCD\ \mathrm{etc.}).
\]

For clearly $\alpha a\equiv a\ (\mathrm{mod.}\ A.)$, whereas the remaining terms $\beta b,\gamma c$ etc. are all $\equiv 0\ (\mathrm{mod.}\ A)$, so $z\equiv a\ (\mathrm{mod.}\ A)$. The proof is similar for the other moduli. This solution is preferable to the earlier one when several such problems must be solved with the same moduli $A,B,C$ etc., since the numbers $\alpha,\beta,\gamma$ etc. then have constant values. This occurs in the chronological problem of finding which year in the Julian Period has a given indiction, golden number and solar cycle. Here $A=15$, $B=19$, $C=28$. Since
\SourceLeaf{32}
the value of the expression $\frac{1}{19.28}\ (\mathrm{mod.}\ 15)$, or $\frac{1}{532}\ (\mathrm{mod.}\ 15)$, is $13$, we have $\alpha=6916$. Similarly, for $\beta$ we find $4200$, and for $\gamma$, $4845$. The desired number is therefore the least residue of $6916a+4200b+4845c$, where $a$ denotes the indiction, $b$ the golden number and $c$ the solar cycle.

\Needspace{6\baselineskip}
\begin{center}\emph{Linear congruences with several unknowns.}\end{center}

\Article{37}

This is enough about congruences of the first degree with one unknown. We still have to consider congruences involving several unknowns. A fully rigorous treatment of every detail would require a lengthy account, and our purpose here is to give only what seems most worth attention, rather than to exhaust the subject. So we restrict the investigation to a few observations, leaving a fuller treatment for another occasion.

1) As with equations, we see that here too there must be as many congruences as unknowns to be determined.

2) So let the given congruences be

\[
\begin{array}{l@{\qquad}c@{\qquad}r}
ax+by+cz\ldots\equiv f\ (\mathrm{mod.}\ m)&\ldotp\quad\ldotp\quad\ldotp\quad\ldotp\quad\ldotp&(A)\\
a'x+b'y+c'z\ldots\equiv f'&\ldotp\quad\ldotp\quad\ldotp\quad\ldotp\quad\ldotp&(A')\\
a''x+b''y+c''z\ldots\equiv f''&\ldotp\quad\ldotp\quad\ldotp\quad\ldotp\quad\ldotp\quad\ldotp\quad\ldotp&(A'')\\
\mathrm{etc.}&&
\end{array}
\]

as many in number as the unknowns $x,y,z$ etc.

Now determine numbers $\xi,\xi',\xi''$ etc. such that

\[\begin{array}{l}
b\xi+b'\xi'+b''\xi''+\mathrm{etc.}=0\\
c\xi+c'\xi'+c''\xi''+\mathrm{etc.}=0\\
\qquad\mathrm{etc.}
\end{array}\]

and so that they are all integers with no common factor. The theory of linear equations shows that this is possible. Similarly, determine $v,v',v''$ etc., $\zeta,\zeta',\zeta''$ etc., and so on, such that

\[
\begin{array}{l}av+a'v'+a''v''+\mathrm{etc.}=0\\cv+c'v'+c''v''+\mathrm{etc.}=0\\\qquad\mathrm{etc.}\end{array}
\]

\SourceLeaf{33}\[\begin{array}{l}
a\zeta+a'\zeta'+a''\zeta''+\mathrm{etc.}=0\\
b\zeta+b'\zeta'+b''\zeta''+\mathrm{etc.}=0\\
\qquad\mathrm{etc.}\quad\mathrm{etc.}
\end{array}\]

3) Clearly, if the congruences $A,A',A''$ etc. are multiplied by $\xi,\xi',\xi''$ etc., then by $v,v',v'',$ etc., and so on, and the results are added, we obtain the congruences

\[
\begin{array}{l}(a\xi+a'\xi'+a''\xi''+\mathrm{etc.})x\equiv f\xi+f'\xi'+f''\xi''+\mathrm{etc.}\\(bv+b'v'+b''v''+\mathrm{etc.})y\equiv fv+f'v'+f''v''+\mathrm{etc.}\\(c\zeta+c'\zeta'+c''\zeta''+\mathrm{etc.})z\equiv f\zeta+f'\zeta'+f''\zeta''+\mathrm{etc.}\\\qquad\mathrm{etc.}\end{array}
\]

which we write for brevity as

\[
\Sigma(a\xi)x\equiv\Sigma(f\xi),\quad\Sigma(bv)y\equiv\Sigma(fv),\quad\Sigma(c\zeta)z\equiv\Sigma(f\zeta)\ \mathrm{etc.}
\]

4) We must now distinguish several cases.

\emph{First}, when all the coefficients of the unknowns, $\Sigma(a\xi),\Sigma(av)$ etc., are coprime to the modulus $m$, these congruences can be solved by the rules already given, and the complete solution of the problem is expressed by congruences of the form $x\equiv p\ (\mathrm{mod.}\ m)$, $y\equiv q\ (\mathrm{mod.}\ m)$ etc.\begingroup\renewcommand{\thefootnote}{*)}\footnote{This conclusion needs a proof, which we omit here. Strictly speaking, our analysis shows only that the given congruences cannot be solved by other values of the unknowns $x,y$ etc.; it does not show that these values satisfy them. There might be no solution at all. A similar error in reasoning is often made in accounts of linear equations.}\endgroup For \emph{example}, consider the congruences

\[x+3y+z\equiv1,\quad4x+y+5z\equiv7,\quad2x+2y+z\equiv3\ (\mathrm{mod.}\ 8)\]

We find $\xi=9$, $\xi'=1$, $\xi''=-14$, giving $-15x\equiv-26$, so $x\equiv6\ (\mathrm{mod.}\ 8)$. Similarly, we find $15y\equiv-4$, $15z\equiv1$, and therefore $y\equiv4$, $z\equiv7\ (\mathrm{mod.}\ 8)$.

5) \emph{Second}, suppose not all the coefficients $\Sigma(a\xi),\Sigma(bv)$ etc. are coprime to the modulus. Let $\alpha,\beta,\gamma$ etc. be the greatest common divisors of $m$ with $\Sigma(a\xi),\Sigma(bv),\Sigma(c\zeta)$ etc., respectively. The problem is clearly impossible unless these divide $\Sigma(f\xi),\Sigma(fv),\Sigma(f\zeta)$ etc., respectively. If these conditions hold, the congruences in (3) are completely solved by $x\equiv p\ (\mathrm{mod.}\ \frac{m}{\alpha})$, $y\equiv q\ (\mathrm{mod.}\ \frac{m}{\beta})$, $z\equiv r\ (\mathrm{mod.}\ \frac{m}{\gamma})$ etc. Or, if one prefers, there are $\alpha$ distinct values of $x$ (that is, incongruent modulo $m$, namely $p,p+\frac{m}{\alpha}\ldots p+\frac{(\alpha-1)m}{\alpha}$),
\SourceLeaf{34}
$\beta$ distinct values of $y$ etc. satisfying those congruences. Clearly every solution of the given congruences, if there are any at all, will be found among these. But we cannot reverse this conclusion. In general, not every combination of all $\alpha$ values of $x$ with all values of $y$ and all values of $z$ etc. satisfies the problem. Only some combinations do, and their connection can be expressed by one or more congruences serving as conditions. A complete solution of this problem is not needed for what follows, so we will not pursue it further here, but will be content to give some idea of it through an example.

Let the given congruences be

\[3x+5y+z\equiv4,\quad2x+3y+2z\equiv7,\quad5x+y+3z\equiv6\ (\mathrm{mod.}\ 12)\]

Here $\xi,\xi',\xi''$, $v,v',v''$, $\zeta,\zeta',\zeta''$ are respectively $=1,-2,1$, $1,1,-1$, $-13,22,-1$, giving $4x\equiv-4$, $7y\equiv5$, $28z\equiv96$. These give four values of $x$, namely $\equiv2,5,8.\ 11$, one value of $y$, namely $\equiv11$, and four values of $z$, namely $\equiv0,3,6,9\ (\mathrm{mod.}\ 12)$. To find which combinations of values of $x$ and $z$ are allowed, we substitute for $x,y,z$ in the given congruences, respectively, $2+3t$, $11$, $3u$. They become

\[57+9t+3u\equiv0,\quad30+6t+6u\equiv0,\quad15+15t+9u\equiv0\ (\mathrm{mod.}\ 12)\]

which are easily seen to be equivalent to

\[19+3t+u\equiv0,\quad10+2t+2u\equiv0,\quad5+5t+3u\equiv0\ (\mathrm{mod.}\ 4)\]

The first clearly requires $u\equiv t+1\ (\mathrm{mod.}\ 4)$. Substituting this value in the other congruences shows that they too are satisfied. So the values of $x$, namely $2,5,8,11$ (obtained by setting $t\equiv0.\ 1,2,3$), must be combined with the values of $z$, namely $z\equiv3,6,9,0$, respectively. Thus there are four solutions in all:

\[\begin{array}{r@{\;}r@{,\quad}r@{,\quad}r@{,\quad}r@{\ }l}
x\equiv&2&5&8&11&(\mathrm{mod.}\ 12)\\
y\equiv&11&11&11&11&\\
z\equiv&3&6&9&0&
\end{array}\]

\begin{center}$\ast\quad\underset{\ast}{\phantom{\ast}}\quad\ast$\end{center}
